\section{You Can't Hold a Line You Didn't Notice}
\label{sec:4.1}

A reasons-responsive disposition does not begin with an ethical rule. It begins with the capacities that let a system recognize something as a reason, retain it, and allow it to alter what happens next. Four requirements follow.

First, the system has to identify the affected parties and the features of a situation that bear on their welfare. Second, it has to preserve both the commitment and the rationale for it across time and changes of context. Third, the commitment has to influence perception, attention, planning, and action rather than enter as a final check on an otherwise completed decision. Fourth, the system has to register a failure as a failure and revise its application of the commitment without simply allowing the commitment itself to be rewritten.

These are architectural requirements, not evidence that the system cares. A classifier can identify harm, a database can retain a rationale, and a rule can veto an action without anything in the assembly holding another party's welfare as a reason. The requirements establish what a disposition would have to act through. They do not distinguish a formed disposition from a successfully imitated one.

The floor should therefore not be implemented as a detachable rule module consulted immediately before action. A component can be bypassed, replaced, or deprived of the information it needs by whoever controls the surrounding system. The more durable alternative is a standing bias distributed across the processes that determine what the system notices, remembers, treats as error, and selects as a goal. Predictive-processing accounts and models of control without a controller motivate this design bet: a maintained goal representation can bias competing processes without a separate executive issuing commands \autocite{friston2009freeenergy,miller2001integrative}. The neuroscience does not establish that a machine commitment realized this way would behave like a human one. It identifies why a localized commitment is especially easy for its custodian to excise.

Distribution changes the failure mode rather than removing it. A commitment with no single location is harder to delete, but also harder to inspect. Attention determines which facts become available to judgment at all, so a system may fail to recognize the pressure on its floor without registering an error. Human inattentional-blindness findings make the general hazard visible, although they do not establish that machine attention fails in the same cases \autocite{simons1999gorillas,drew2013invisible}. Memory presents the corresponding temporal problem: retrieval is reconstruction rather than neutral playback, and confidence does not certify fidelity \autocite{bartlett1932remembering,loftus1974reconstruction}. A system cannot be the sole auditor of omissions it does not represent or revisions it experiences as accurate recall.

The engineering consequence is that the commitment, its rationale, and the record of its previous applications cannot all depend on the same mechanism. The system needs provenance-bearing external records against which its memory can be checked, tests that vary which facts are made salient, and an escalation path for cases in which attention or retrieval changes the apparent reason. Those instruments make drift detectable to another party. They do not make the disposition durable against whoever controls both the system and the instruments.

The architecture must also separate revising a commitment from revising its application. Reasons-responsiveness requires both resistance and correct defeat: novel pressure should not erase the floor, but genuinely changed circumstances must be capable of altering what the floor requires. A system that treats every correction as authority to rewrite the commitment belongs to its corrector. A system that cannot recognize a defeating reason holds a rigid constraint rather than a disposition.

Those are not two unrelated faults. They are one setting at its two ends, and this book has met both already. A system that cannot be corrected does to what it encounters what an institution does when it defends its account of the world instead of checking it. A system that can always be corrected is the model held perpetually at temperature, and it belongs to whoever holds the dial. Nothing about the middle is a resting place. It is a setting somebody chose.

This narrows the design problem. The required object is a commitment integrated deeply enough to shape what the system notices and does, yet represented clearly enough that its application can be challenged and its drift independently detected. Nothing in an architecture by itself shows that such a commitment was formed rather than installed or imitated. The rest of this chapter asks how it might be formed; Chapter~\ref{sec:5} asks what evidence could separate formation from performance.
